% =====================================================================
\chapter{Cuestiones}
\label{chap:parteA}
% =====================================================================

Cada cuestión vale 1~punto y tiene una sola respuesta correcta. Las cuestiones A.1 a A.5 corresponden al módulo~12; A.6 a A.10, al módulo~13; A.11 a A.18, al módulo~14, y A.19 y A.20, al módulo~15.

\begin{enumerate}[label=\textbf{A.\arabic*.},ref=A.\arabic*,leftmargin=2.6em,itemsep=8pt]

% ---------------------------------------------------------------- M12
\item\label{q:bpe-paso} ¿Qué hace el algoritmo BPE en cada paso de su entrenamiento? \punto
\begin{enumerate}[label=(\alph*)]
  \item Elimina del vocabulario el token menos probable.
  \item Divide la palabra más larga del corpus en dos mitades.
  \item Fusiona el par de símbolos adyacentes más frecuente en el corpus.
  \item Asigna a cada palabra un entero según su frecuencia.
\end{enumerate}

\item\label{q:bytes} Sin contar los tokens especiales, un tokenizador BPE a nivel de byte parte de un vocabulario de: \punto
\begin{enumerate}[label=(\alph*)]
  \item 128 símbolos.
  \item 256 símbolos.
  \item 65\,536 símbolos.
  \item tantos símbolos como caracteres distintos aparezcan en el corpus.
\end{enumerate}

\item\label{q:pinon} ¿Cuántos bytes ocupa la palabra «piñón» codificada en UTF-8? \punto
\begin{enumerate}[label=(\alph*)]
  \item 5.
  \item 6.
  \item 7.
  \item 10.
\end{enumerate}

\item\label{q:orden} Al codificar una palabra nueva con un tokenizador BPE ya entrenado, las reglas de fusión se aplican: \punto
\begin{enumerate}[label=(\alph*)]
  \item en el orden en que se aprendieron.
  \item de la que produce el token más largo a la que produce el más corto.
  \item en orden aleatorio, para dar variedad a la segmentación.
  \item solo la regla más frecuente.
\end{enumerate}

\item\label{q:rnn} ¿Por qué las redes recurrentes escalan peor que el Transformer? \punto
\begin{enumerate}[label=(\alph*)]
  \item Porque tienen más parámetros por capa.
  \item Porque no pueden usar embeddings.
  \item Porque solo admiten vocabularios de caracteres.
  \item Porque procesan los tokens uno tras otro, lo que impide paralelizar el entrenamiento a lo largo de la secuencia, y sus gradientes se desvanecen con la distancia.
\end{enumerate}

% ---------------------------------------------------------------- M13
\item\label{q:forma} En la atención de una cabeza sobre $n$ tokens, la matriz $A=\softmax\bigl(QK\T/\sqrt{d_k}\bigr)$ tiene forma: \punto
\begin{enumerate}[label=(\alph*)]
  \item $n\times d_k$.
  \item $d_k\times d_k$.
  \item $n\times n$.
  \item $n\times d_v$.
\end{enumerate}

\item\label{q:varianza} Si las componentes de $q,k\in\R^{d_k}$ son variables independientes de media~0 y varianza~1, la varianza de $q\cdot k$ es: \punto
\begin{enumerate}[label=(\alph*)]
  \item $1$.
  \item $d_k$.
  \item $\sqrt{d_k}$.
  \item $d_k^{2}$.
\end{enumerate}

\item\label{q:mascara} La máscara causal se implementa: \punto
\begin{enumerate}[label=(\alph*)]
  \item poniendo a~0 las puntuaciones de los tokens futuros antes de la softmax.
  \item poniendo a~0 los pesos de los tokens futuros después de la softmax, sin renormalizar.
  \item eliminando de $V$ las filas de los tokens futuros.
  \item sumando $-\infty$ a las puntuaciones de los tokens futuros antes de la softmax.
\end{enumerate}

\item\label{q:mha} Con $d=768$ y 12 cabezas, ¿cuántos parámetros tiene la atención multi-cabeza, sin sesgos? \punto
\begin{enumerate}[label=(\alph*)]
  \item 196\,608.
  \item 589\,824.
  \item 1\,179\,648.
  \item 2\,359\,296.
\end{enumerate}

\item\label{q:contexto} Pasar de un contexto de 1024 a uno de 8192 tokens multiplica la memoria de la matriz de atención, almacenada completa, por: \punto
\begin{enumerate}[label=(\alph*)]
  \item 8.
  \item 64.
  \item 512.
  \item 4096.
\end{enumerate}

% ---------------------------------------------------------------- M14
\item\label{q:prenorm} En un bloque pre-norm, la subcapa de atención se escribe: \punto
\begin{enumerate}[label=(\alph*)]
  \item $h=x+\Attn(\Norm(x))$.
  \item $h=\Norm(x+\Attn(x))$.
  \item $h=\Norm(\Attn(x))$.
  \item $h=\Attn(x)+\Norm(x)$.
\end{enumerate}

\item\label{q:normfinal} ¿Por qué un Transformer pre-norm necesita una normalización final antes del unembedding? \punto
\begin{enumerate}[label=(\alph*)]
  \item Porque la norma del flujo residual crece con la profundidad.
  \item Porque la softmax exige entradas positivas.
  \item Para poder atar los pesos de entrada y de salida.
  \item No la necesita; solo la llevan los modelos post-norm.
\end{enumerate}

\item\label{q:swiglu} Un MLP SwiGLU tiene el mismo número de parámetros que un MLP clásico de dimensión oculta $4d$ cuando su dimensión oculta vale: \punto
\begin{enumerate}[label=(\alph*)]
  \item $2d$.
  \item $\tfrac{8}{3}d$.
  \item $3d$.
  \item $4d$.
\end{enumerate}

\item\label{q:rope-rel} RoPE consigue que la puntuación $\tilde q_m\cdot\tilde k_n$ entre el token~$m$ y el token~$n$: \punto
\begin{enumerate}[label=(\alph*)]
  \item dependa de las posiciones solo a través de $n-m$.
  \item dependa de las posiciones solo a través de $m+n$.
  \item sea independiente de las posiciones.
  \item crezca linealmente con la posición.
\end{enumerate}

\item\label{q:rope-v} RoPE gira las consultas y las claves, pero no los valores, porque: \punto
\begin{enumerate}[label=(\alph*)]
  \item $v$ tiene dimensión impar.
  \item girar $v$ cambiaría su norma.
  \item la posición debe influir en a quién se atiende, no en qué contenido se copia.
  \item $v$ no participa en la atención.
\end{enumerate}

\item\label{q:llama-emb} LLaMA-2~7B no ata sus embeddings ($V=32\,000$, $d=4096$). ¿Cuántos parámetros suman la matriz de entrada y la de salida? \punto
\begin{enumerate}[label=(\alph*)]
  \item 65\,536\,000.
  \item 131\,072\,000.
  \item 262\,144\,000.
  \item 524\,288\,000.
\end{enumerate}

\item\label{q:16bytes} Entrenar con Adam en precisión mixta (pesos y gradientes en bf16, copia maestra en fp32 y momentos $m$ y $v$ en fp32) ocupa, sin contar las activaciones: \punto
\begin{enumerate}[label=(\alph*)]
  \item 4 bytes por parámetro.
  \item 8 bytes por parámetro.
  \item 12 bytes por parámetro.
  \item 16 bytes por parámetro.
\end{enumerate}

\item\label{q:6nd} Según la regla $6ND$, ¿cuántos FLOPs cuesta, aproximadamente, entrenar un modelo de $10^{9}$ parámetros con $2\cdot 10^{10}$ tokens? \punto
\begin{enumerate}[label=(\alph*)]
  \item $6\cdot 10^{18}$.
  \item $2\cdot 10^{19}$.
  \item $4\cdot 10^{19}$.
  \item $1{,}2\cdot 10^{20}$.
\end{enumerate}

% ---------------------------------------------------------------- M15
\item\label{q:lnV} Un GPT con $V=50\,257$ tokens, recién inicializado con logits casi nulos, debe mostrar una pérdida inicial (en nats por token) cercana a: \punto
\begin{enumerate}[label=(\alph*)]
  \item $10{,}82$.
  \item $1$.
  \item $0$.
  \item $50\,257$.
\end{enumerate}

\item\label{q:checkpoint} Durante el entrenamiento, la pérdida de entrenamiento sigue bajando y la de validación empieza a subir. ¿Qué punto de control debe conservarse como modelo final? \punto
\begin{enumerate}[label=(\alph*)]
  \item El último.
  \item El de menor pérdida de validación.
  \item El de menor pérdida de entrenamiento.
  \item El primero, porque aún no ha sobreajustado.
\end{enumerate}

\end{enumerate}
